\section{Compute-Optimal Test-Time Scaling}

\subsection{Muestreo paralelo y revisión secuencial}

Aumentar el test-time compute tiene dos formas básicas:

\begin{itemize}
    \item \textbf{Parallel sampling}: genera varios candidatos de manera independiente a partir del mismo prompt, con énfasis en la diversidad;
    \item \textbf{Sequential revision}: hace que el modelo lea los resultados intermedios y los revise, con énfasis en la profundidad de una sola trayectoria.
\end{itemize}

\videofigure{figures/parallel-sampling.jpg}{El muestreo paralelo usa varios candidatos independientes para ampliar el ancho de búsqueda.}{00:27:00--00:29:06}

Los métodos paralelos son fáciles de escalar y los errores no se propagan entre candidatos; los métodos secuenciales pueden aprovechar la retroalimentación para mejorar de forma progresiva, pero si la dirección inicial es errónea, los pasos posteriores pueden seguir optimizando alrededor de una premisa equivocada.

\subsection{Outcome Reward Model y Process Reward Model}

\videofigure{figures/reward-model-types.jpg}{Outcome reward y process reward ofrecen señales de búsqueda con distinta granularidad.}{00:29:06--00:30:08}

\begin{itemize}
    \item \textbf{ORM}（Outcome Reward Model）: solo evalúa la respuesta final; es adecuado para best-of-$N$, con asignación de crédito burda;
    \item \textbf{PRM}（Process Reward Model）: evalúa el proceso de razonamiento paso a paso; puede usarse en beam search y para podar de forma anticipada.
\end{itemize}

Si una trayectoria contiene los pasos $z_{1:T}$, el PRM puede ofrecer una puntuación local $r_t=r(z_t\mid z_{<t},x)$, con la cual el buscador conserva las ramas de mayor potencial, sin necesidad de esperar a que se genere la respuesta completa para descubrir el error.

\subsection{Combinar ancho y profundidad}

Una estrategia compute-optimal no consiste en elegir de forma fija entre paralelo o secuencial, sino en combinar ambos dentro de un presupuesto total: primero explorar varios candidatos iniciales y luego someter a los candidatos de mayor puntuación a varias rondas de revisión.

\videofigure{figures/hybrid-search.jpg}{El muestreo paralelo amplía el ancho, y el revision model aumenta la profundidad de la búsqueda.}{00:31:10--00:35:28}

Puede escribirse como una restricción de presupuesto:

$$
B=n_{\mathrm{parallel}}\cdot c_{\mathrm{sample}}
+n_{\mathrm{revise}}\cdot c_{\mathrm{revise}}
+c_{\mathrm{verify}}.
$$

\begin{itemize}
    \item $B$: presupuesto total de inferencia para un problema;
    \item $n_{\mathrm{parallel}}$: número de candidatos paralelos;
    \item $n_{\mathrm{revise}}$: número de rondas de revisión secuencial;
    \item $c_{\mathrm{sample}},c_{\mathrm{revise}},c_{\mathrm{verify}}$: el costo de cada operación.
\end{itemize}

\subsection{Conclusiones experimentales: la mejor estrategia varía según la dificultad}

\videofigure{figures/compute-optimal-results.jpg}{Guiado por el PRM, el compute-optimal search ofrece un beneficio inicial mayor que el best-of-$N$ puro.}{00:35:28--00:37:28}

\videofigure{figures/parallel-sequential-ratio.jpg}{Los problemas sencillos se inclinan hacia la revisión secuencial, mientras que los difíciles requieren cierta exploración paralela.}{00:37:28--00:40:12}

La intuición es la siguiente:

\begin{itemize}
    \item los problemas sencillos suelen tener ya una solución inicial cercana a la correcta, y la revision puede corregirla rápidamente;
    \item la solución inicial de los problemas difíciles puede caer en un planteamiento totalmente equivocado, por lo que se necesita explorar en paralelo distintas direcciones;
    \item un exceso de secuencialidad desperdicia el presupuesto en reparar trayectorias irrecuperables;
    \item un exceso de paralelismo produce una gran cantidad de candidatos superficiales, sin aprovechar suficientemente la retroalimentación.
\end{itemize}

\subsection{Cuándo los FLOPs de test-time superan a los FLOPs de pre-training}

\videofigure{figures/testtime-vs-pretraining.jpg}{Los problemas de dificultad media suelen beneficiarse del test-time compute, mientras que los problemas más difíciles siguen limitados por la capacidad de base.}{00:40:12--00:42:32}

\begin{importantbox}{El test-time scaling no sustituye por completo al pre-training scaling}
Cuando el modelo base casi no tiene trayectorias correctas para cierto tipo de problema, la búsqueda no puede amplificar una señal que no existe. Los experimentos muestran que, para los problemas más difíciles, el beneficio de aumentar el cómputo de inferencia es muy pequeño, y ampliar el preentrenamiento o mejorar los datos de entrenamiento sigue siendo más eficaz.
\end{importantbox}

\subsection{Resumen de la sección}

La estrategia de inferencia óptima depende de la dificultad del problema, la capacidad del modelo y la calidad del verifier. Un sistema verdaderamente compute-optimal necesita estimar en línea el beneficio marginal y decidir de forma dinámica si continúa muestreando, revisando, verificando o se detiene.
